\chapter{A Concrete Naming Certificate for $\CauchyNetDiagonalUp$}
\label{ch:concrete-instances-cauchy-net-diagonal-namecert}
\origin{human}

Following Moore--Smith and Kelley on nets, the $\CauchyNetDiagonalUp$ packet names the finite diagonal extraction route by which a displayed Cauchy net is read through a cofinal schedule and handed to regular rational names. It sits over the Cauchy-net source of \autoref{ch:concrete-instances-cauchynetcompletion-namecert}, the Moore--Smith net surface of \autoref{ch:concrete-instances-moore-smith-completion-namecert}, the stream-window naming surface of \autoref{ch:concrete-instances-streamname-namecert}, the regular rational readback of \autoref{ch:concrete-instances-regseqrat-namecert}, the dyadic tolerance scale of \autoref{ch:concrete-instances-dyadic-tolerance-scale-namecert}, and the Real seal of \autoref{ch:concrete-instances-real-namecert}. The chapter records a finite diagonal packet only; it does not add a choice functional, quotient net equality, selected subnet, selected limit, or ambient $\RealUp$ completeness theorem.

\begin{definition}[Cauchy-net diagonal carrier]
\label{def:cauchy-net-diagonal-carrier}
A $\CauchyNetDiagonalUp$ carrier is a finite $\BHist$ packet
$$
D=(A,M,K,T,S,Q,R,H,C,P,N)
$$
with the following displayed rows:
\begin{enumerate}
\item $A$ is the $\CauchyNetUp$ source row whose directed-domain and Cauchy evidence are displayed before any diagonal consumer reads the packet.
\item $M$ is the $\MooreSmithNetUp$ source row, recording that the source is consumed as Moore--Smith net data rather than as an unordered family.
\item $K$ is the finite cofinal diagonal schedule, listing the observed directed indices and the finite cofinality checks used by the packet.
\item $T$ is the dyadic tolerance ledger, recording the requested error levels and the finite comparisons used to meet them.
\item $S$ is the $\StreamNameUp$ window row obtained by reading the scheduled observations in order.
\item $Q$ is the $\RegSeqRatUp$ readback attached to those stream windows.
\item $R$ is the $\RealUp$ seal row reached after the regular rational readback.
\item $H$ records componentwise $\hsame$ transport for $A,M,K,T,S,Q,R,C,P,N$.
\item $C$ records displayed $\Cont$ replay from the Cauchy-net source through the Moore--Smith row, finite schedule, dyadic ledger, stream window, regular rational readback, and Real seal.
\item $P$ is the $\Pkg$ provenance over $\CauchyNetUp$, $\MooreSmithNetUp$, $\DyadicRatCoreUp$, $\StreamNameUp$, $\RegSeqRatUp$, $\RealUp$, $\BHist$, $\hsame$, $\Cont$, and $\NameCert$.
\item $N$ is the local $\NameCert_{\CauchyNetDiagonalUp}$ row.
\end{enumerate}
The classifier compares accepted packets only through $A,M,K,T,S,Q,R,H,C,P,N$. It has no coordinate for a choice functional, quotient net equality, selected subnet, selected representative, selected limit, quotient stream equality, host real equality, or ambient $\RealUp$ completeness.
\end{definition}
\leandef{BEDC.Derived.CauchyNetDiagonalUp.CauchyNetDiagonalUp}

\begin{theorem}[Cauchy-net diagonal NameCert obligations]
\label{thm:cauchy-net-diagonal-namecert-obligations}
For every accepted $\CauchyNetDiagonalUp$ carrier $D=(A,M,K,T,S,Q,R,H,C,P,N)$, the local $\NameCert_{\CauchyNetDiagonalUp}$ surface consists exactly of the following obligation rows:
\begin{enumerate}
\item carrier admission for the displayed Cauchy-net source, Moore--Smith source, finite cofinal diagonal schedule, dyadic tolerance ledger, stream window, regular rational readback, Real seal, transport, replay, provenance, and local naming rows;
\item source visibility, requiring the rows $A$ and $M$ to remain available before any diagonal schedule or completion-facing read is consumed;
\item schedule exactness, requiring $K$ to display the finite cofinal observations and $T$ to display the dyadic tolerance ledger used by those observations;
\item readback routing, requiring $S,Q,R$ to be consumed in order from stream windows to regular rational data and then to the Real seal;
\item structural stability, requiring $H,C,P,N$ to preserve and name exactly the same finite route.
\end{enumerate}
No NameCert obligation is stored outside these rows.
\end{theorem}
\begin{proof}
Carrier admission is projection from \autoref{def:cauchy-net-diagonal-carrier}. The accepted packet exposes exactly the rows $A,M,K,T,S,Q,R,H,C,P,N$, so the local naming surface can read only those coordinates.

The source rows come first. The Cauchy-net row $A$ supplies directed Cauchy evidence, while the Moore--Smith row $M$ fixes the net interface through which the diagonal route is read.

The finite diagonal step is then local to $K$ and $T$. The schedule lists the observed directed indices, and the dyadic ledger lists the tolerance comparisons that justify each observed window.

Finally the packet reads $S,Q,R$ as stream windows, regular rational data, and a Real seal. Componentwise $\hsame$, displayed $\Cont$ replay, $\Pkg$ provenance, and the local $\NameCert$ row preserve this same route without adding any hidden selector or completion theorem.
\end{proof}
\leanchecked{BEDC.Derived.CauchyNetDiagonalUp.CauchyNetDiagonalNameCertObligations}

\begin{theorem}[Cauchy-net diagonal RegSeqRat handoff]
\label{thm:cauchy-net-diagonal-regseqrat-handoff}
Every regular-rational or Real-facing consumer of an accepted $\CauchyNetDiagonalUp$ carrier reads the diagonal route only through
$$
\CauchyNetUp,\MooreSmithNetUp,\DyadicRatCoreUp,\StreamNameUp,\RegSeqRatUp,\RealUp,\BHist,\hsame,\Cont,\Pkg,\NameCert .
$$
The handoff exports no choice functional, quotient net equality, selected subnet, selected representative, selected limit, quotient stream equality, host real equality, or ambient $\RealUp$ completeness.
\end{theorem}
\begin{proof}
Start with the accepted carrier. The only source coordinates are the Cauchy-net row $A$ and the Moore--Smith row $M$, so a consumer cannot read a detached family, selected subnet, or quotient net before the displayed source route is fixed.

The diagonal observations pass through $K$ and $T$. These rows are finite data: the cofinal observation schedule and the dyadic tolerance ledger. They do not contain a choice operator or a global completion principle.

The handoff to regular rational names is exactly the route $S,Q,R$. The stream row supplies the windows, the regular rational row reads those windows, and the Real row seals the resulting name.

The remaining rows $H,C,P,N$ transport, replay, source, and name the same packet. Since none of these rows is a selector, quotient equality, host equality, or ambient completeness theorem, the RegSeqRat and Real-facing handoff is exhausted by the displayed finite route.
\end{proof}

\closureat{CauchyNetDiagonalUp}{seedStr}

\begin{closurestatus}{\CauchyNetDiagonalUp}
  \theoryclosure{\seedClosure}
  \scopeclosed{\autoref{def:cauchy-net-diagonal-carrier}, \autoref{thm:cauchy-net-diagonal-namecert-obligations}, and \autoref{thm:cauchy-net-diagonal-regseqrat-handoff} seed the finite diagonal extraction route from Cauchy-net and Moore--Smith data to StreamName, RegSeqRat, and Real readback.}
  \formalstatus{\theoremCheckedV}
  \leantarget{BEDC.Derived.CauchyNetDiagonalUp.CauchyNetDiagonalNameCertObligations}
  \bridgestatus{none}
  \notclaimed{This chapter does not close a choice functional, quotient net equality, selected subnet, selected representative, selected limit, quotient stream equality, host real equality, ambient $\RealUp$ completeness, Lean verification, or bridge checking.}
  \upgradepath{Obligation closure requires checked carrier exposure, source visibility, finite cofinal schedule exactness, dyadic tolerance ledger exactness, readback routing, transport, replay, provenance, and local naming for BEDC.Derived.CauchyNetDiagonalUp.}
  \constructivestory{A $\CauchyNetDiagonalUp$ packet is generated from a CauchyNet row, a MooreSmithNet source row, a finite cofinal diagonal schedule, a dyadic tolerance ledger, StreamName windows, RegSeqRat readback, a Real seal, componentwise hsame transport, displayed Cont replay, Pkg provenance, and a local NameCert row.}
\end{closurestatus}
